\chapter{What money rewards, and what it cannot see}\label{ch:money}

A useful account must enter a decision somewhere. This chapter follows a workshop choosing between a paid order and an essential repair. The example separates three things that will remain distinct throughout the book: what can physically happen, how its consequences are evaluated, and which rule makes someone willing or able to act.

\section{The decision at the workshop}
The workshop can repair the town's water line today or complete an order from a customer who will pay immediately. Its workers need wages, and its tools and electricity cost money. The paid order is a serious claim on its time. So is the pipe: households and a clinic depend on it.

A simplified private calculation is
\begin{equation}
 \text{expected private return}
 =\text{expected receipts}-\text{private expenditure}.
\end{equation}
This expression leaves out many reasons people work, including care, duty, curiosity and legal obligation. It nevertheless matters to an organisation that must meet its bills. A repair may protect many people while leaving the workshop unable to cover its own expenditure. Another activity may pay well while shifting physical costs onto people outside the agreement.

Economics describes an effect on others that a market price does not include as an externality. Pollution is a familiar example. The IMF's explanations of externalities and energy subsidies discuss circumstances in which environmental costs are not reflected in the price paid by a user.\cite{intro-externality,intro-imf-fossil} The omitted consequence remains a physical consequence whether or not an invoice names it.

For EBU, the practical starting point is to declare the boundary before evaluating the change. A repair that restores a receiver while depleting an unrecorded source has an incomplete physical description. A fuller account can expose the omitted effect. Turning that information into a contract, payment or obligation is a subsequent institutional decision.

\section{A dollar is a claim, not a physical forecast}
Money makes claims transferable. A deposit can pay a supplier who then pays someone else, without each payment carrying a complete description of the earlier activity. Modern bank lending can create deposits, subject to the constraints described by the Bank of England.\cite{intro-boe} Those monetary mechanisms do not themselves forecast the stock of safe water or working equipment available when the money is spent.

Imagine two orders bringing the same revenue to the workshop. One uses a scarce spare part in a short-lived installation; the other repairs a pump needed for continuing service. The financial inflow is equal, but the future physical possibilities differ. A good contract or public rule may take that difference into account. It needs a description of the difference before it can do so reliably.

Nor does observing a favourable condition prove how it was produced. Water may arrive because of today's repair, yesterday's dispatch or an external change upstream. Later chapters will distinguish the physical outcome from a response model and from causal attribution. That distinction begins here: a receipt, a forecast and an observation are different kinds of record.

A central EBU feature is therefore traceability. A claim about improvement should point to the state, boundary and completed transition on which it rests. This can make duplicated claims and missing consequences easier to detect. It may reduce the cost of verification, but it does not remove the need for trusted measurements or accountability when a record is wrong.

\section{Demand, desire and indispensable service}
Offers in a market depend on purchasing power as well as preference. An additional convenience for one person can attract a larger offer than an indispensable service for another. This does not establish what producers feel about either person. It identifies a limitation of using offers alone to rank needs.

Public provision, insurance, subsidies, rights and regulation can alter that allocation. Their existence shows that an economy already contains many mechanisms for connecting payment to broader purposes. The question is how those mechanisms learn whether the promised service was actually delivered and whether maintaining it has displaced an equally important need elsewhere.

Consider a clinic whose water is restored for an hour and then fails again. A final stock reading may look satisfactory at the chosen observation time, while the interruption has still mattered to patients. The eventual EBU endpoint account and a measure of service continuity would answer different questions. A responsible institution may need both. Naming the distinction early prevents a single physical number from acquiring a role its definition does not support.

\section{Where an incentive points}
Suppose a supplier earns ten monetary units by exporting water and six by supplying the clinic. If the clinic's consequences do not enter the supplier's contract, export has the larger private return. A service obligation can change that choice. A drought can change which withdrawals are physically possible. These are three separate changes: in a payment comparison, an institutional rule and a physical constraint.

The proposed account can inform such decisions by describing a declared physical transition consistently. It cannot manufacture water, appoint an owner or decide whose need takes priority. Those boundaries make the account more useful: the same physical result can be inspected even when people disagree about the appropriate rule.

\begin{workedbox}{Two ledgers for one decision}
A repair costs eight monetary units and earns five today, giving private return $-3$. If observations also show that it prevents a specified water loss, that fact does not change the arithmetic of the financial account. It gives the community evidence for designing a payment, duty or public service that makes the repair possible.
\end{workedbox}

The efficiency opportunity is to connect incentives to verifiable outcomes rather than merely to activity labels. A rule that rewards reported repairs needs a way to distinguish completed restoration, repeated reporting of the same event and damage followed by repair. The later laws of finite valuation and sequential closure will provide part of that checking structure. The institution can then make its maintenance and access rules explicit, using the physical history as their common reference.

\section{The next difficulty: who can reach what exists?}
Even an accurate account of quantity and price does not establish access. A resource can be present but unusable, unreachable, already committed or withheld under existing rights. The next chapter separates these failures so that the response addresses the actual obstacle.
